\begin{pastexam}[label={px:ncu-109-dsa-10}]{中央資工所 109 學年度 第 10 題}
（是非；答對 1 分，答錯倒扣 0.5 分）

Depth-first search can be implemented by a stack. \qquad (A)~True. \qquad (B)~False.
\end{pastexam}

\begin{pastsoln}
答案：(A) True。

遞迴版的 DFS 本來就是靠呼叫堆疊：每次往下拜訪一個點就推入一層，那個點的鄰點都看完了就彈出，回到
上一個點繼續。 把呼叫堆疊換成自己的堆疊，就是資結書第~21.4.6~節的迴圈版（虛擬碼~21.4）：堆疊裡
放「目前路徑上的點，以及每個點的鄰接串列看到哪裡」，每次看堆疊頂端的點，還有沒拜訪的鄰點就推入，
沒有就彈出。 這樣拜訪的順序和遞迴版完全相同。

對照一下：把 DFS 的堆疊換成佇列，得到的是 BFS（資結書第~21.3~節）。 這正是兩者的差別：DFS 每次從
最後放進去的點繼續往下走，BFS 每次處理最早放進去的點。
\end{pastsoln}
